\documentclass[11pt,a4paper]{article}
% Společný preambule pro cvičnou písemku a klíč řešení (XeLaTeX)
\usepackage[a4paper,left=2.1cm,right=2.1cm,top=2.4cm,bottom=2.2cm,headheight=14pt,footskip=1.1cm]{geometry}
\usepackage{amsmath}
\usepackage{fontspec}
\usepackage{unicode-math}
\usepackage{polyglossia}
\setdefaultlanguage{czech}
\setmainfont{Latin Modern Roman}
\setmathfont{Latin Modern Math}
\setsansfont{Inter}[Scale=0.92]
\usepackage{microtype}
\usepackage{xcolor}
\usepackage{enumitem}
\usepackage{array}
\usepackage{multicol}
\usepackage{needspace}
\usepackage{xspace}
\usepackage{fancyhdr}
\usepackage{lastpage}
\usepackage{tikz}
\usetikzlibrary{calc}
\usepackage[skins,breakable]{tcolorbox}

\definecolor{ink}{HTML}{1F2937}
\definecolor{accent}{HTML}{2B4C7E}
\definecolor{accentlight}{HTML}{EAF0FA}
\definecolor{muted}{HTML}{6B7280}
\definecolor{rulegray}{HTML}{B8C2D0}
\definecolor{rulelight}{HTML}{E1E6EE}
\definecolor{good}{HTML}{1F6F4A}
\definecolor{goodlight}{HTML}{EAF6EF}
\definecolor{warnlight}{HTML}{FFF6E5}
\definecolor{warn}{HTML}{9A5B00}

\color{ink}
\linespread{1.12}
\setlength{\parindent}{0pt}
\setlength{\parskip}{0pt}
\raggedright
\setlist{nosep}

% Záhlaví a zápatí
\pagestyle{fancy}
\fancyhf{}
\renewcommand{\headrulewidth}{0pt}
\renewcommand{\footrulewidth}{0pt}
\fancyhead[L]{\sffamily\footnotesize\color{muted}Matematické dovednosti · Množiny}
\fancyfoot[C]{\sffamily\footnotesize\color{muted}\thepage\,/\,\pageref{LastPage}}

% Pruh s názvem části
\newcommand{\cast}[2]{%
  \par\Needspace*{11cm}\vspace{1.6em}\noindent
  \begin{tcolorbox}[enhanced,colback=accent,colframe=accent,boxrule=0pt,arc=3pt,left=10pt,right=10pt,top=5pt,bottom=5pt,before skip=0pt,after skip=6pt]
    {\sffamily\bfseries\color{white}\small Část #1\enspace\textcolor{white!70!accent}{|}\enspace #2}
  \end{tcolorbox}\par\nopagebreak}

% Úloha: číslo, text
\newcommand{\uloha}[1]{\par\vspace{1.15em}\noindent\textcolor{accent}{\sffamily\bfseries #1}\enspace}
% Body
\newcommand{\bodu}[1]{\hfill\mbox{\sffamily\footnotesize\color{muted}#1\,b}}

% Místo na řešení: linkovaný box dané výšky
\newcommand{\misto}[1]{%
  \par\vspace{0.35em}
  \begin{tcolorbox}[enhanced,height=#1,colback=white,colframe=rulegray,boxrule=0.5pt,arc=3pt,left=0pt,right=0pt,top=0pt,bottom=0pt,before skip=0pt,after skip=0pt,
    underlay={\begin{scope}\clip (frame.south west) rectangle (frame.north east);
      \foreach \i in {1,...,40}{\draw[rulelight,line width=0.4pt] ([yshift=-\i*0.9cm]frame.north west) -- ([yshift=-\i*0.9cm]frame.north east);}
    \end{scope}}]
  \end{tcolorbox}\par}

% Rámeček s poznámkou
\newtcolorbox{tip}[1][]{enhanced,breakable,colback=accentlight,colframe=accentlight,boxrule=0pt,arc=4pt,left=10pt,right=10pt,top=7pt,bottom=7pt,before skip=8pt,after skip=8pt,#1}
\newtcolorbox{varovani}[1][]{enhanced,breakable,colback=warnlight,colframe=warnlight,boxrule=0pt,arc=4pt,left=10pt,right=10pt,top=7pt,bottom=7pt,before skip=8pt,after skip=8pt,#1}

% Řešení (pro klíč)
\newtcolorbox{reseni}{enhanced,breakable,colback=goodlight,colframe=goodlight,boxrule=0pt,arc=3pt,left=10pt,right=8pt,top=6pt,bottom=6pt,before skip=5pt,after skip=2pt,
  borderline west={2.2pt}{0pt}{good}}
\newcommand{\pravda}{\textcolor{good}{\sffamily\bfseries pravda}\xspace}
\newcommand{\nepravda}{\textcolor{warn}{\sffamily\bfseries nepravda}\xspace}

% Zkratky
\newcommand{\N}{\mathbb N}
\newcommand{\Z}{\mathbb Z}
\newcommand{\Pot}{\mathcal P}
\newcommand{\pr}{\mathrel{\mid}}
\newcommand{\nepr}{\mathrel{\nmid}}
\newcommand{\sm}{\setminus}

% Krátké položky ve dvou sloupcích / jeden sloupec
\newenvironment{dvasl}{\begin{multicols}{2}\begin{enumerate}[label=(\alph*),leftmargin=2.2em,itemsep=3pt,topsep=3pt]}{\end{enumerate}\end{multicols}}
\newenvironment{pol}{\begin{enumerate}[label=(\alph*),leftmargin=2.2em,itemsep=3pt,topsep=4pt]}{\end{enumerate}}
\setlength{\columnsep}{1.2cm}

\newcommand{\blacksquare}{\mathord{\rule{0.85ex}{0.85ex}}}

\fancyhead[R]{\sffamily\footnotesize\color{muted}Cvičná písemka}

\begin{document}

% ---------- Titulní blok ----------
{\sffamily\bfseries\fontsize{25}{28}\selectfont\color{accent} Množiny}\par
\vspace{2pt}
{\sffamily\large\color{muted} Cvičná písemka s místem na řešení}\par
\vspace{10pt}
{\color{rulegray}\rule{\linewidth}{0.6pt}}\par
\vspace{10pt}

\begin{tip}
{\sffamily\bfseries\small Jak s tím pracovat}\par\vspace{3pt}
\begin{itemize}[leftmargin=1.2em,itemsep=2pt]
  \item \textbf{Části 1–5} procvičují jednotlivá témata (zápisy množin, prvky a $\emptyset$, podmnožiny, sjednocení a průniky, potenční množina). Piš přímo do linkovaných polí.
  \item \textbf{Část 6} je cvičná písemka na čas: \textbf{20 minut, 20 bodů, bez pomůcek}. Dělej ji až nakonec.
  \item Řešení jsou ve druhém souboru \textit{Cvicna-pisemka-mnoziny-reseni.pdf}. Nahlédni do nich až po vyřešení.
  \item Použité značení: $\N=\{1,2,3,\dots\}$, $\N_0=\N\cup\{0\}$, $x\pr y$ znamená „$x$ dělí $y$“.
\end{itemize}
\end{tip}

% ==================== ČÁST 1 ====================
\cast{1}{Zápisy množin}

\Needspace*{7.6cm}
\uloha{1.1}\textbf{Zapište výčtem prvků:}
\begin{pol}
  \item $A=\{x\in\Z\mid x^2<20\}$
  \item $B=\{n\in\N\mid n\pr 24\ \land\ n>3\}$
  \item $C=\{3k-1\mid k\in\{1,2,3,4\}\}$
  \item $D=\{x\in\N\mid 3\pr x\ \land\ x<20\}$
\end{pol}
\misto{4.2cm}

\Needspace*{9.7cm}
\uloha{1.2}\textbf{Zapište množinovým zápisem} (u (a) a (b) alespoň dvěma způsoby):
\begin{pol}
  \item množinu všech násobků pěti mezi přirozenými čísly,
  \item množinu všech druhých mocnin přirozených čísel,
  \item množinu všech přirozených čísel, která dávají po dělení třemi zbytek 2,
  \item množinu všech přirozených dělitelů čísla 30.
\end{pol}
\misto{6.3cm}

\Needspace*{10.6cm}
\uloha{1.3}\textbf{Nechť $X,Y\subseteq\N$. Zapište množinovým zápisem:}
\begin{pol}
  \item množinu čísel z $X$, která jsou dělitelná \emph{všemi} čísly z $Y$,
  \item množinu čísel z $Y$, která jsou menší než \emph{některé} číslo z $X$,
  \item množinu čísel z $X$, která jsou větší než \emph{každé} číslo z $Y$,
  \item množinu všech součtů dvou (ne nutně různých) čísel z $X$,
  \item množinu čísel z $Y$, která jsou druhou mocninou nějakého čísla z $X$.
\end{pol}
\misto{7.2cm}

\Needspace*{7.0cm}
\uloha{1.4}Prvek leží v \emph{právě jedné} z množin $A$, $B$. Zapište množinu všech takových prvků (i)~pomocí operací $\cup,\cap,\setminus$, (ii)~pomocí $\in$, $\notin$ a logických spojek.
\misto{3.6cm}

\Needspace*{6.6cm}
\uloha{1.5}Popište slovy, jaká čísla obsahuje množina
\[M=\{n\in\N\mid \forall k\in\N:\ (k\pr n\Rightarrow k=1\lor k=n)\}.\]
(Pozor na nejmenší prvek.)
\misto{3.2cm}

% ==================== ČÁST 2 ====================
\cast{2}{Prvky, prázdná množina, počet prvků}

\Needspace*{8.8cm}
\uloha{2.1}\textbf{Rozhodněte, zda je výrok pravdivý, a krátce zdůvodněte:}
\begin{dvasl}
  \item $\emptyset\subseteq\{\emptyset\}$
  \item $\{\emptyset\}\subseteq\emptyset$
  \item $\emptyset\in\{\{\emptyset\}\}$
  \item $\{\emptyset\}\in\{\{\emptyset\}\}$
  \item $\{1,2\}\in\{1,2,\{1,2\}\}$
  \item $\{1,2\}\subseteq\{1,2,\{1,2\}\}$
  \item $\{\{1\}\}\subseteq\{1,\{1\}\}$
  \item $\Pot(\emptyset)=\{\emptyset\}$
\end{dvasl}
\misto{5.4cm}

\Needspace*{8.8cm}
\uloha{2.2}\textbf{Kolik prvků mají tyto množiny? Prvky vypište.}
\begin{pol}
  \item $\{1,\{1\},\{1,1\},\{\{1\}\},\emptyset,\{\emptyset\},\{\}\}$
  \item $\{a,b\}\cup\{b,c\}\cup\{\{a,b\}\}$
  \item $\{\emptyset,\{\emptyset,\emptyset\},\{\{\emptyset\}\}\}$
  \item $\Pot(\{a,b\})\cup\{a,b\}$
\end{pol}
\misto{5.4cm}

\Needspace*{9.0cm}
\uloha{2.3}Nechť $S=\{1,\,2,\,\{3\},\,\{1,2\}\}$. Do tabulky napište \emph{ano} / \emph{ne}:
\par\vspace{5pt}
\begin{center}
\renewcommand{\arraystretch}{2.15}
\setlength{\tabcolsep}{14pt}
\begin{tabular}{@{}|>{\centering\arraybackslash}p{2.6cm}|>{\centering\arraybackslash}p{2.6cm}|>{\centering\arraybackslash}p{2.6cm}|@{}}
\hline
{\sffamily\small\color{muted}$x$} & {\sffamily\small\color{muted}$x\in S$} & {\sffamily\small\color{muted}$x\subseteq S$}\\
\hline
$\{1\}$ & & \\ \hline
$\{3\}$ & & \\ \hline
$\{1,2\}$ & & \\ \hline
$\{\{3\}\}$ & & \\ \hline
$\emptyset$ & & \\ \hline
\end{tabular}
\end{center}

\Needspace*{8.4cm}
\uloha{2.4}\begin{pol}
  \item Vypište všechny prvky množiny $\Pot(\{1,\{2\}\})$.
  \item Vypište $\Pot(\Pot(\emptyset))$.
  \item Kolik prvků má $\Pot(\Pot(\{a\}))$?
\end{pol}
\misto{5cm}

% ==================== ČÁST 3 ====================
\cast{3}{Podmnožiny}

\Needspace*{6.7cm}
\uloha{3.1}Nechť $A=\{1,2,3,4,5\}$, $B=\{2,4\}$, $C=\{1,2,3\}$, $D=\{2,3,6\}$. Rozhodněte pravda/nepravda:
\begin{dvasl}
  \item $B\subseteq A$
  \item $D\subseteq A$
  \item $B\subseteq C$
  \item $C\cap D\subseteq C$
  \item $C\cap D\subseteq B$
  \item $A\setminus D\subseteq A$
\end{dvasl}
\misto{3.3cm}

\Needspace*{8.4cm}
\uloha{3.2}\textbf{Zapište pomocí kvantifikátorů} (nad prvky množin $A$, $B$):
\begin{dvasl}
  \item $A\subseteq B$
  \item $A\nsubseteq B$
  \item $A\cap B=\emptyset$
  \item $A\subsetneq B$
  \item $A\neq B$
\end{dvasl}
\misto{5cm}

\Needspace*{8.0cm}
\uloha{3.3}Nechť $A=\{n\in\N\mid 6\pr n\}$, $B=\{n\in\N\mid 2\pr n\}$, $C=\{n\in\N\mid 3\pr n\}$. Určete vztah (platí $\subseteq$, $\supseteq$, $=$, nebo jsou množiny neporovnatelné?) a zdůvodněte:
\begin{pol}
  \item $A$ a $B$ \qquad (b) $A$ a $C$ \qquad (c) $B$ a $C$ \qquad (d) $A$ a $B\cap C$
\end{pol}
\misto{4.6cm}

\Needspace*{11.0cm}
\uloha{3.4}\textbf{Dokažte:}
\begin{pol}
  \item Jestliže $A\subseteq B$, pak $A\cap C\subseteq B\cap C$.
  \item Jestliže $A\cup B=A\cap B$, pak $A=B$.
\end{pol}
\misto{7.6cm}

% ==================== ČÁST 4 ====================
\cast{4}{Sjednocení, průnik, rozdíl}

\Needspace*{8.0cm}
\uloha{4.1}Nechť $A=\{1,2,3,4,5,6\}$, $B=\{2,4,6,8\}$, $C=\{3,6,9\}$. Vypište:
\begin{dvasl}
  \item $A\cup B$
  \item $A\cap B$
  \item $A\setminus B$
  \item $B\setminus A$
  \item $(A\cap B)\cup C$
  \item $A\setminus(B\cup C)$
  \item $(A\cup C)\setminus B$
  \item $|A\cup B\cup C|$
\end{dvasl}
\misto{4.6cm}

\Needspace*{7.6cm}
\uloha{4.2}Nechť $A=\{n\in\N\mid 2\pr n\}$, $B=\{n\in\N\mid 3\pr n\}$, $C=\{n\in\N\mid n\le12\}$. Vypište prvky:
\begin{pol}
  \item $A\cap B\cap C$ \quad (b) $(A\cup B)\cap C$ \quad (c) $C\setminus(A\cup B)$ \quad (d) $(A\setminus B)\cap C$
\end{pol}
\misto{4.2cm}

\Needspace*{11.8cm}
\uloha{4.3}Které z množin jsou si pro libovolné $A,B,C$ rovny? Rozdělte je do skupin. Jednu rovnost dokažte (řetězem ekvivalencí) a u jedné dvojice, která si rovna není, uveďte konkrétní protipříklad.
\begin{dvasl}
  \item $A\setminus(B\cap C)$
  \item $(A\setminus B)\cup(A\setminus C)$
  \item $(A\setminus B)\cap(A\setminus C)$
  \item $A\setminus(B\cup C)$
  \item $(A\cup B)\setminus C$
  \item $(A\setminus C)\cup(B\setminus C)$
\end{dvasl}
\misto{8.4cm}

\Needspace*{11.6cm}
\uloha{4.4}Pro libovolné množiny $A,B,C$: dokažte, nebo vyvraťte protipříkladem.
\begin{dvasl}
  \item $(A\setminus B)\cup B=A$
  \item $A\cap(B\setminus C)=(A\cap B)\setminus C$
  \item $A\setminus(A\setminus B)=A\cap B$
  \item $A\setminus(B\setminus C)=(A\setminus B)\setminus C$
\end{dvasl}
\misto{8.2cm}

\Needspace*{8.2cm}
\uloha{4.5}\begin{pol}
  \item Je-li $|A|=12$, $|B|=9$ a $|A\cap B|=4$, určete $|A\cup B|$, $|A\setminus B|$, $|A\times B|$ a $|\Pot(A\cap B)|$.
  \item Ve třídě je 30 studentů. Fotbal hraje 18, basketbal 14 a oba sporty 7 studentů. Kolik jich nehraje ani jeden?
\end{pol}
\misto{4.8cm}

% ==================== ČÁST 5 ====================
\cast{5}{Potenční množina, indexovaná sjednocení a průniky}

\Needspace*{7.0cm}
\uloha{5.1}Nechť $A=\{1,2\}$. Rozhodněte pravda/nepravda:
\begin{dvasl}
  \item $\{1\}\in\Pot(A)$
  \item $1\in\Pot(A)$
  \item $\emptyset\in\Pot(A)$
  \item $\{\{1\}\}\subseteq\Pot(A)$
  \item $A\subseteq\Pot(A)$
  \item $\{1,2,3\}\in\Pot(A)$
\end{dvasl}
\misto{3.6cm}

\Needspace*{9.8cm}
\uloha{5.2}Pro $i=1,2,\dots,5$ nechť $M_i=\{j\in\N\mid 2j\le i\}$. Vypište $M_1,\dots,M_5$ a určete:
\begin{pol}
  \item $\displaystyle\bigcup_{i=1}^{5}M_i$ \qquad (b) $\displaystyle\bigcap_{i=1}^{5}M_i$ \qquad (c) $\displaystyle\bigcup_{i=1}^{5}\{M_i\}$ \qquad (d) $\displaystyle\bigcap_{i=1}^{5}\{M_i\}$ \qquad (e) $\displaystyle\Big\{\bigcup_{i=1}^{5}M_i\Big\}$
\end{pol}
\misto{6.4cm}

\Needspace*{11.0cm}
\uloha{5.3}Pro libovolné množiny $A,B$:
\begin{pol}
  \item dokažte $\Pot(A\cap B)=\Pot(A)\cap\Pot(B)$,
  \item na příkladu ukažte, že $\Pot(A\cup B)\neq\Pot(A)\cup\Pot(B)$ nemusí platit rovnost,
  \item platí vždy $\Pot(A)\cup\Pot(B)\subseteq\Pot(A\cup B)$? Zdůvodněte.
\end{pol}
\misto{7.6cm}

\Needspace*{7.8cm}
\uloha{5.4}\textbf{Kolik prvků mají:}
\begin{pol}
  \item $\Pot(\{1,2,3,4,5\})$,
  \item množina všech podmnožin $\{1,2,3,4,5\}$ se sudým počtem prvků,
  \item množina všech dvouprvkových podmnožin $\{1,2,3,4,5\}$,
  \item $\Pot(A\times B)$, kde $|A|=2$ a $|B|=3$?
\end{pol}
\misto{4.4cm}

% ==================== ČÁST 6 ====================
\newpage
\cast{6}{Cvičná písemka na čas}

\begin{tcolorbox}[enhanced,colback=white,colframe=accent,boxrule=0.8pt,arc=4pt,left=12pt,right=12pt,top=7pt,bottom=7pt,before skip=2pt,after skip=4pt]
\sffamily\small
\begin{tabular}{@{}p{0.31\linewidth}p{0.31\linewidth}p{0.3\linewidth}@{}}
\textbf{Čas:} 20 minut & \textbf{Max.:} 20 bodů & \textbf{Bez pomůcek} \\
\end{tabular}
\par\vspace{4pt}
{\color{muted}Zapište celé postupy a zdůvodnění. U kvantifikátorů vždy uveďte, přes jakou množinu kvantifikujete.}
\end{tcolorbox}
\begin{center}
\sffamily\small
\renewcommand{\arraystretch}{1.6}
\setlength{\tabcolsep}{10pt}
\begin{tabular}{|c|c|c|c|c|c|}
\hline
\color{muted}Úloha 1 & \color{muted}Úloha 2 & \color{muted}Úloha 3 & \color{muted}Úloha 4 & \color{muted}Úloha 5 & \textbf{Celkem}\\
\hline
\hspace{1cm}/\,4 & \hspace{1cm}/\,4 & \hspace{1cm}/\,4 & \hspace{1cm}/\,4 & \hspace{1cm}/\,4 & \hspace{1.2cm}/\,20\\
\hline
\end{tabular}
\end{center}


\Needspace*{7.6cm}
\uloha{Úloha 1}\textbf{Zápis množin.}\bodu{4}\par
Zapište množinovým zápisem ($X,Y\subseteq\N$, $M$ je konečná množina):
\begin{pol}
  \item množinu všech přirozených čísel dělitelných čtyřmi \emph{nebo} šesti,
  \item množinu čísel z $X$, která jsou menší než každé číslo z $Y$,
  \item množinu všech trojnásobků čísel z $X$,
  \item množinu všech podmnožin $M$, které mají alespoň dva prvky.
\end{pol}
\misto{4.2cm}

\Needspace*{7.0cm}
\uloha{Úloha 2}\textbf{Pravda, nebo nepravda?}\bodu{4}\par
Zdůvodněte krátce:
\begin{pol}
  \item $\{\emptyset\}\subseteq\Pot(\{1\})$ \qquad (b) $\emptyset\in\{\{\emptyset\}\}$ \qquad (c) $\{\{1\}\}\subseteq\{1,\{1\}\}$ \qquad (d) $2\in\Pot(\{1,2\})$
\end{pol}
\misto{3.6cm}

\Needspace*{7.8cm}
\uloha{Úloha 3}\textbf{Počty prvků.}\bodu{4}\par
\begin{pol}
  \item Kolik prvků má $\{a,\{a\},\{a,a\},\emptyset,\{\emptyset\},\{\{a\}\}\}$? Vypište je. \hfill{\sffamily\footnotesize\color{muted}2 b}
  \item $|\Pot(\Pot(\{1\}))|=\ $? \hfill{\sffamily\footnotesize\color{muted}1 b}
  \item Pro $A=\{1,2,3\}$, $B=\{3,4\}$ určete $|A\times B|$ a $|A\cup B|$. \hfill{\sffamily\footnotesize\color{muted}1 b}
\end{pol}
\misto{4.4cm}

\Needspace*{7.6cm}
\uloha{Úloha 4}\textbf{Kvantifikátory a negace.}\bodu{4}\par
Zapište formulí s kvantifikátory větu „Každé číslo z množiny $A$ je menší než některé číslo z množiny $B$“ ($A,B\subseteq\N$) a poté formuli \textbf{znegujte} (bez negace před kvantifikátorem). \hfill{\sffamily\footnotesize\color{muted}2 b + 2 b}
\misto{4.2cm}

\Needspace*{11.0cm}
\uloha{Úloha 5}\textbf{Rovnosti množin.}\bodu{4}\par
\begin{pol}
  \item Dokažte: $(A\cup B)\setminus C=(A\setminus C)\cup(B\setminus C)$. \hfill{\sffamily\footnotesize\color{muted}2 b}
  \item Vyvraťte: $A\setminus(B\setminus C)=(A\setminus B)\setminus C$. \hfill{\sffamily\footnotesize\color{muted}2 b}
\end{pol}
\misto{7.6cm}



\end{document}
